% !Mode:: "TeX:UTF-8"
% !Mode:: "TeX:UTF-8"
\newcommand{\thepage}{} % to silence a hyperref warning about pdfpagelabels
\RequirePackage[]{silence}
\WarningsOff[hyperref]
\PassOptionsToPackage{log-declarations=false}{xparse}
%\RequirePackage[l2tabu, orthodox]{nag}
\RequirePackage{luatex85} % to make \pdf...= commands work with LuaTeX
\documentclass[aspectratio=169,t]{beamer}
\usepackage{ifthen}
\usepackage[T1]{fontenc}
\usepackage{iftex}
\ifLuaTeX
 \usepackage{fontspec}
 \setmainfont[Ligatures=TeX]{Source Sans Pro}
\else
 \usepackage[utf8]{inputenc}
\fi
\usepackage{german}
\usepackage{newunicodechar}
\newunicodechar{€}{\texteuro}
\newunicodechar{ß}{\ss}
\renewcommand{\SS}{\iffontchar\font"1E9E \symbol{"1E9E}\else SS\fi}
\usepackage{ellipsis}
\usepackage{microtype}
\usepackage{listings}
\usepackage{sourcecodepro}
\lstset{language=Python, basicstyle=\ttfamily}
\usepackage{amsfonts}
\usepackage{amssymb}
\usepackage{amsthm}
\usepackage{fancyhdr}
\usepackage{graphicx}

\newtheorem{proposition}[theorem]{Proposition}

\DeclareMathOperator*{\argmax}{argmax}
\DeclareMathOperator{\supp}{supp}
\newcommand{\E}{\ensuremath{\mathbb{E}}}
\newcommand{\V}{\ensuremath{\mathbb{V}}}

\setbeamercolor{emph}{fg=purple}
\renewcommand<>{\emph}[1]{%
  {\usebeamercolor[fg]{emph}\only#2{\itshape}#1}%
}



\hypersetup{draft=false,colorlinks=false,pdfpagemode=UseOutlines,plainpages=false,
            breaklinks=true,bookmarksnumbered=true,pdfborder=0 0 0,
            pdftitle={Data Science SoSe 2026 HHU},
            pdfauthor={Konrad Voelkel},
            pdfsubject={data science}}

%\documentclass[aspectratio=169,t]{beamer}
\usetheme{hhubeamer} % put the hhu-beamer-theme in /usr/share/texmf/tex/latex/beamer/themes/

% -------------------- Code listings config --------------------
% Define custom colors for syntax highlighting
\definecolor{codegreen}{rgb}{0,0.6,0}
\definecolor{codegray}{rgb}{0.5,0.5,0.5}
\definecolor{codepurple}{rgb}{0.58,0,0.82}
\definecolor{backcolour}{rgb}{0.96,0.96,0.94} % Very light gray/beige
% Configure the listings package
\lstset{
    language=Python,
    backgroundcolor=\color{backcolour},   
    commentstyle=\color{codegreen}\itshape,
    keywordstyle=\color{blue}\bfseries,
    numberstyle=\tiny\color{codegray},
    stringstyle=\color{codepurple},
    basicstyle=\ttfamily\footnotesize, % Use a monospaced font
    breakatwhitespace=false,         
    breaklines=true,      
             % Enable automatic line breaking
    captionpos=b,                      % Caption at the bottom
    keepspaces=true,                 
    numbers=left,                      % Line numbers on the left
    numbersep=10pt,                    % Distance of line numbers from the code
    showspaces=false,                
    showstringspaces=false,            % Don't show weird underscores in strings
    showtabs=false,                  
    tabsize=4,
    frame=single,                      % Add a subtle border around the code
    rulecolor=\color{black!20},        % Border color
  }

\event[]{}
\author[Sommersemester 2026]{Dr. Konrad Völkel, Sommersemester 2026}

%\title{Data Science, VL 1}
%\date{16.04.2026}


%%% Local Variables:
%%% mode: latex
%%% TeX-master: t
%%% End:


\title{Data Science, VL 15}
\subtitle{Schätztheorie}
\date{16.06.2026}

\begin{document}

\begin{frame}
  \titlepage
\end{frame}

\section{Statistische Modelle}

\begin{frame}[separator]
  \frametitle{Problem first:\\ verschiedene Schätzungen}
\end{frame}


\begin{frame}
  \frametitle{Einstieg: Standardabweichung schätzen}

Gegeben eine Stichprobe $x$ von $n$ \texttt{float}s, die aus einer univariaten normalverteilten Zufallsvariable gezogen sind, wollen wir die Parameter dieser Normalverteilung schätzen:
{\Large
  $\mu(X)$ und $\sigma^2(X)$.
}

\uncover<-1>{
  {\tiny Beispiel
    $x = \texttt{np.array}$
    
\hspace{2.4cm}$\begin{matrix}
[ 6.5281 & 3.8003 & 4.9575 & 7.4818 & 6.7351 & 1.0454 \\
4.9002 & 2.6973 & 2.7936 & 3.8212 & 3.2881 & 5.9085 \\
4.5221 & 3.2434 & 3.8877 & 3.6673 & 5.9882 & 2.5897 \\
3.6261 & 1.2918 & -2.1060 & 4.3072 & 4.7289 & 1.5157 \\
7.5395 & 0.0913 & 3.0915 & 2.6256 & 6.0656 & 5.9387 \\
3.3099 & 3.7563 & 1.2244 & -0.9616 & 2.3042 & 3.3127 \\
5.4606 & 5.4048 & 2.2253 & 2.3954 & 0.9029 & 0.1600 \\
-0.4125 & 6.9016 & 1.9807 & 2.1239 & 0.4944 & 4.5550 ]
\end{matrix}$
}
}
\uncover<2->{\vspace{-2.2cm}}

\uncover<2->{Geschenkt: $\mu(X)$ schätzen wir mit $\texttt{mean}(x)$.}

\uncover<3>{(Wieso?)}

\uncover<4->{
    Es gibt zwei übliche Formeln, um $\sigma^2(X)$ aus $x$ zu schätzen:
  \[\sigma(x)^2 = \dfrac{1}{n} \sum_{i=1}^n (x_i-\mu)^2,\qquad \sigma^\ast(x)^2 = \dfrac{1}{n-1} \sum_{i=1}^n (x_i-\mu)^2\]
}

\uncover<5->{ Moment mal: was ist $\mu$? }
\uncover<6->{ Entweder $\mu(X)$ oder $\texttt{mean}(x)$.}
\end{frame}

\begin{frame}
  \frametitle{Einstieg: Standardabweichung schätzen}
  Gegeben eine Stichprobe $x$ von $n$ \texttt{float}s, die aus einer univariaten normalverteilten Zufallsvariable gezogen sind, wollen wir die Parameter dieser Normalverteilung schätzen:
{\Large
  $\mu(X)$ und $\sigma^2(X)$.
}

\vspace{0.5cm}

\uncover<1->{Wenn $x = (x_1,\dots,x_n)$ gegeben ist, können wir $\texttt{mean}(x)$ berechnen. Umgekehrt genügen $n-1$ Werte $(x_1 - \texttt{mean}(x), \dots, x_{n-1} - \texttt{mean}(x))$ zur Berechnung von $x_n - \texttt{mean}(x) = \sum_{i=1}^{n-1}(\texttt{mean}(x) - x_i)$}


\vspace{0.5cm}

\uncover<2->{Beobachtung: gegeben $(x_1,\dots,x_{n-1})$ sind $\texttt{mean}(x)$ und $x_n$ gleich viel 'wert'}.
\end{frame}

\begin{frame}
  \frametitle{Einstieg: Standardabweichung schätzen}
  Gegeben eine Stichprobe $x$ von $n$ \texttt{float}s, die aus einer univariaten normalverteilten Zufallsvariable gezogen sind, wollen wir die Parameter dieser Normalverteilung schätzen:
{\Large
  $\mu(X)$ und $\sigma^2(X)$.
}

\vspace{0.5cm}

\uncover<2>{Wenn $\mu = \mu(X)$ bekannt ist, dann ist 
  \[\sigma(x)^2 = \dfrac{1}{n} \sum_{i=1}^n (x_i-\mu)^2\]
  eine Schätzung für $\sigma^2(X)$, die für größere Stichproben immer besser wird.}

\uncover<3->{
\vspace{-2.4cm}

  Wenn $\mu = \texttt{mean}(x)$ geschätzt ist, dann ist 
  \[\sigma(x)^2 = \dfrac{1}{n} \sum_{i=1}^n (x_i-\mu)^2\]
  eine Schätzung für $\sigma^2(X)$, die für größere Stichproben immer besser wird.}

\uncover<4->{Aber: der Erwartungswert von $\sigma(x)^2$ für zufällige Stichproben ist nicht $\sigma^2(X)$,
  sondern $\dfrac{n-1}{n}\sigma^2(X)$. Anders bei $\mu = \mu(X)$, dann klappt es!
  }
\end{frame}

\begin{frame}
  \frametitle{Einstieg: Standardabweichung schätzen}
  Wir können den Schätzer korrigieren (``Bessel-Korrektur''):
  \[\sigma^\ast(x)^2 = \dfrac{n}{n-1} \sigma(x)^2\]

  \uncover<2->{
    Wir sprechen bei $\sigma(x)^2$ von einem \alert{asymptotisch erwartungstreuen Schätzer} für $\sigma(X)^2$. Der Schätzer $\sigma^\ast(x)^2$ ist sogar \alert{erwartungstreu} (frei von \alert{Bias}).
    
    [Das besprechen wir noch ausführlicher]
  }

\vspace{1cm}
  
  \uncover<3->{
    Wir besprechen jetzt einmal den mathematischen Formalismus, wie er auch in Scikit-Learn zur Anwendung kommt.
    }
  \end{frame}

\begin{frame}
  \frametitle{Was ist Schätzen?}

  \begin{itemize}[<+->]
  \item  Geschätzt wird eine Größe (wie $\mu(X)$) aus einer Stichprobe (wie $x$), um aus einer Familie von Modellen (z.B. Normalverteilungen) eine bestimmte bzw. eine Teilfamilie auszuwählen (die mit dem Parameter $\mu = \mu(x)$).
  \item Wir formalisieren die 'Größe' als Abbildung, die Familie als \alert{Statistisches Modell} und die (zufällige) Stichprobe als (recht 'große') Zufallsvariable.
  \item  Das 'Schätzen' ist eine Prozedur, die zu jeder möglichen Stichprobe einen Wert ausspuckt. Auch das modellieren wir als Zufallsvariable!
  \end{itemize}
\end{frame}  

  \begin{frame}
    \frametitle{Statistisches Modell}
    \begin{definition}
      Ein statistisches Modell ist ein Tripel $(\mathcal{X}, \mathcal{F}, P_\theta : \theta \in \Theta)$
      aus
      \begin{itemize}
        \item
          einem Stichprobenraum $\mathcal{X}$,
        \item einer $\sigma$-Algebra $\mathcal{F}$ auf $\mathcal{X}$ und
        \item einer Klasse $\{P_\theta\}_{\theta \in \Theta}$ von Wahrscheinlichkeitsmaßen auf $(\mathcal{X}, \mathcal{F})$,\\ die mindestens zwei Elemente enthält, also $|\Theta| > 1$.
        \end{itemize}
        Wir schreiben oft einfach $(\mathcal{X}, P_\theta : \theta \in \Theta)$.

    \end{definition}
  \end{frame}

  \begin{frame}
    \frametitle{Parametrische Modelle, Standardmodelle}
    \begin{definition}
      Ein statistisches Modell $(\mathcal{X}, P_\theta : \theta \in \Theta)$ heißt \alert{parametrisches Modell}, wenn die Parametermenge $\Theta$ eine Teilmenge eines $\mathbb{R}^d$ ist für ein $d \in \mathbb{N}$. Für $d=1$ sagt man auch \alert{einparametriges Modell}.
    \end{definition}

\uncover<3>{
    \begin{definition}    
Wenn die Stichprobenmenge $\mathcal{X}$ diskret ist (höchstens abzählbar), und $\mathcal{F} = \mathcal{P}(X)$, nennt man das Modell \alert{diskret} und in dem Fall haben alle $P_\theta$ eine pmf $p_\theta$, d.h. $p_\theta(x) = P_\theta(\{x\})$.
\end{definition}
}

\vspace{-2.5cm}

\uncover<4>{
    \begin{definition}
Wenn die Stichprobenmenge $\mathcal{X}$ eine (Borel)-meßbare Teilmenge des $\mathbb{R}^n$ ist (z.B. eine offene Menge oder die Niveaumenge einer stetigen Funktion) und $\mathcal{F}$ eine Einschränkung der Borel-$\sigma$-Algebra von $\mathbb{R}^n$ ist, und außerdem jedes $P_\theta$ eine pdf $f_\theta$ besitzt, dann heißt das statistische Modell ein \alert{stetiges Modell}.
\end{definition}
}

\vspace{-3cm}

\uncover<5>{
    \begin{definition}
Ein statistisches Modell $(\mathcal{X}, P_\theta : \theta \in \Theta)$, welches entweder diskret oder stetig ist (es also pmf oder pdf gibt), nennt man auch \alert{Standardmodell}.
\end{definition}
}
\end{frame}

    \begin{frame}
    \frametitle{Produktmodell}
    Das mehrfache Wiederholen eines Experiments als eigentlich-ein-Experiment umetikettieren und eine Abfolge von Einzelergebnissen dann ein Ergebnis nennen:

    Zu jedem statistischen Modell lässt sich das \alert{Produktmodell} konstruieren:

$$
(\mathcal{X}^N, \mathcal{F}^{\otimes N}, P_{\theta}^{\otimes N} : \theta \in \Theta).
$$

\uncover<2->{
Die Projektionen $X_i \colon \mathcal{X}^N \to \mathcal{X}$ sind dann unabhängig identisch verteilte Zufallsvariablen bezüglich jeder Verteilung $P_\theta$.

$P_\theta^{\otimes N}$ ordnet einem Ereignis in $\mathcal{X}^N$ von der Form $A_1 \times \cdots \times A_N$ das Produkt $\prod_{i=1}^N P_\theta(A_i)$ zu. Damit ist die $N$-fache Wiederholung modelliert.
}
  \end{frame}

    \begin{frame}
      \frametitle{Beispiel}
      \begin{itemize}
        \item  Wir modellieren einen Würfelwurf eines potentiell unfairen Würfels mit einem diskreten Modell: $\mathcal{X} = \{1,\dots,6\}$ und $\Theta = \Delta_6$, die Menge aller möglichen diskreten Verteilungen auf einer sechselementigen Menge (geschrieben als Wahrscheinlichkeitsvektoren in $\mathbb{R}^6$), und $P_\theta(k) = \theta_k$.
        \item  Konkret für $\theta = (\theta_1,\dots,\theta_n)$ ist also die Wahrscheinlichkeit, eine $6$ zu werfen im Modell mit Parameter $\theta$ gegeben durch $P_\theta(6) = \theta_6$.
      \item  Den zweifach wiederholten Würfelwurf modellieren wir also mit $\Theta = \Delta_6$ und $\mathcal{X}^2$, der Menge aller geordneten Paare aus $\mathcal{X}$.

        Dabei ist $P_\theta^{\otimes 2}(5,6) = P_\theta(5)P_\theta(6) = \theta_5\theta_6$.
      \end{itemize}
    \end{frame}
  
    \begin{frame}
    \frametitle{Beispiel}
$\mathcal{X}=\mathbb{R}^N$ mit $\Theta = \mathbb{R} \times \mathbb{R}_{> 0}$ und $P_{\theta}^{\otimes N} = P_{\mu,\sigma}^{\otimes N}$ mit $P_{\mu,\sigma}$ eine univariate Normalverteilung ist das Standardmodell einer Stichprobe der Größe $N$ einer univariaten Zufallsvariable, die normalverteilt ist.
\end{frame}

    \begin{frame}
    \frametitle{Statistik - als mathematischer Begriff}
\begin{definition}
  Sei $(\mathcal{X}, P_\theta : \theta \in \Theta)$ ein statistisches Modell und $\Sigma$ ein weiterer Ereignisraum (d.h. eine Ergebnismenge).

  Dann heißt jede beliebige Zufallsvariable $S \colon \mathcal{X} \to \Sigma$ eine \alert{Statistik} des statistischen Modells.
\end{definition}
\end{frame}


    \begin{frame}
    \frametitle{Schätzer - als mathematischer Begriff}
\begin{definition}
  Sei $(\mathcal{X}, P_\theta : \theta \in \Theta)$ ein statistisches Modell und $\Sigma$ ein weiterer Ereignisraum.
  
  Sei $\tau \colon \Theta \to \Sigma$ eine beliebige Abbildung, die jedem Parameter $\theta \in \Theta$ eine Kenngröße $\tau(\theta) \in \Sigma$ zuordnet.
  
Dann nennen wir jede Statistik $T \colon \mathcal{X} \to \Sigma$ einen \alert{Schätzer} für $\tau$  (auch \alert{Punktschätzer}).
\end{definition}
\end{frame}

      \begin{frame}
        \frametitle{Beispiel}
        \begin{itemize}
          \item
Im Standardmodell einer univariaten Normalverteilung mit Stichprobengröße $N$ ist die Projektion $\Theta \to \mathbb{R}$ auf die erste Komponente, also den Erwartungswert $\mu$ eine Abbildung $\tau$, deren Schätzer uns interessieren.
\item
  Ein möglicher Schätzer ist die Statistik $c \colon \mathbb{R}^N \to \mathbb{R}$, die konstant $0$ ist.
\item
Ein anderer möglicher Schätzer ist das arithmetische Mittel $\texttt{mean} \colon \mathbb{R}^N \to \mathbb{R}$, für den tatsächlich gilt, dass $\texttt{mean}(x)$ für große $N$ beliebig nah an $\mu(X)$ herankommt, wenn $X$ die gesampelte normalverteilte Zufallsvariable ist.
\end{itemize}
  \end{frame}


      \begin{frame}
    \frametitle{Beispiel}
    In den meisten statistischen Modellen können wir die Abbildung $\sigma \colon \Theta \to \mathbb{R}_{\geq 0}$ betrachten, die einem Parameter $\theta$ die Standardabweichung $\sigma(P_\theta)$ der Verteilung $P_\theta$ zuordnet.

    Dafür kennen wir die beiden Schätzer $\sigma,\ \sigma^\ast \colon \mathcal{X} \to \mathbb{R}_{\geq 0}$.

    \vspace{1cm}
    (Erinnerung Bessel-Korrektur: $\sigma^\ast = \sqrt{\frac{n}{n-1}} \sigma$).
  \end{frame}


      \begin{frame}
    \frametitle{Bias und Erwartungstreue}
\begin{definition}
Ein Schätzer $T \colon \mathcal{X} \to \Sigma$ für $\tau \colon \Theta \to \Sigma$ heißt \alert{erwartungstreu} wenn $\mathbb{E}_{\theta}( T ) = \tau(\theta)$ für alle $\theta \in \Theta$ gilt. Gilt das nicht, heißt der Schätzer \alert{verzerrt} (biased) und die Differenz ist der \alert{Bias} $\mathbb{B}_\theta(T) := \mathbb{E}_{\theta}( T ) - \tau(\theta)$.
\end{definition}

\begin{definition}
Der \alert{mittlere quadratische Fehler} des Schätzers zum zu schätzenden Wert ist

$$
\texttt{mse}(T,\theta) := \mathbb{E}_\theta\left ( (T-\tau(\theta))^2 \right ) = \mathbb{V}_\theta(T) + (\mathbb{B}_\theta(T))^2
$$
\end{definition}


  \end{frame}


      \begin{frame}
\frametitle{Folgen von Schätzern}
\begin{definition}
Eine Folge von Schätzern $T_n \colon \mathcal{X}^n \to \Sigma$ für $\tau$ in $n$-fachen Produktmodellen heißt \alert{konsistent} wenn für $n\to\infty$
der Schätzer stochastisch gegen $\tau$ konvergiert: $\lim_{n\to\infty} T_n \xrightarrow{P_\theta} \tau(\theta)$.

Stochastische Konvergenz heißt $\forall \epsilon > 0 : \lim_{n\to\infty} P(|T_n - \tau(\theta)| > \epsilon) = 0$.
\end{definition}

\begin{definition}
Die Folge $T_n$ heißt \alert{asymptotisch erwartungstreu}, wenn $\lim_{n\to\infty} \mathbb{E}_{\theta}(T_n) = \tau(\theta)$.
\end{definition}
  \end{frame}


      \begin{frame}
    \frametitle{Beispiel}
Im Standardmodell einer $N$-fach wiederholten univariaten Normalverteilung ist das Stichprobenmittel $\texttt{mean}(X_1,\dots,X_N)$ ein erwartungstreuer Schätzer für den Parameter $\mu$.
  \end{frame}


      \begin{frame}
        \frametitle{Beispiel}
        \begin{itemize}
\item    Für die Varianz ist der Schätzer $\sigma_N^2 = \frac{1}{N} \sum_{i=1}^N (X_i - \texttt{mean}(X_1,\dots,X_N))^2$ nicht erwartungstreu, denn $\mathbb{E}(\sigma_N^2) = \frac{N-1}{N}\sigma^2(X)$.
\item   Der Bias ist also $-\frac{\sigma^2(X)}{n}$, was für $n\to\infty$ verschwindet. Damit ist der Schätzer asymptotisch erwartungstreu.
\item Der `korrigierte' Schätzer $\frac{N}{N-1} \sigma_N^2 = \frac{1}{N-1} \sum_{i=1}^N (X_i - \texttt{mean}(X_1,\dots,X_N))^2$ ist erwartungstreu.
\end{itemize}
\end{frame}


      \begin{frame}
    \frametitle{Gleichmäßig bester erwartungstreuer Schätzer}
\begin{definition}
Ein zu einer zu schätzenden Funktion $\tau \colon \Theta \to \mathbb{R}$ erwartungstreuer Schätzer $S$ mit endlicher Varianz heißt \alert{gleichmäßig bester erwartungstreuer Schätzer} für $\tau$, wenn für jeden weiteren erwartungstreuen Schätzer $T$ für $\tau$ mit endlicher Varianz für alle $\theta \in \Theta$ gilt:
$\mathbb{V}_\theta(S) \leq \mathbb{V}_\theta(T)$
\end{definition}

Ein solcher Schätzer hat also minimale Streuung um den zu schätzenden Wert unter allen erwartungstreuen Schätzern.
  \end{frame}


  
      \begin{frame}
    \frametitle{Ausblick: Wie finden wir den?}
    Mit der relativen Entropie berechnet man die Fisher-Information und die Cramer-Rao-Ungleichung für reguläre statistische Modelle gibt eine Varianzschranke.

    Für sogenannte exponentielle Familien lässt sich damit direkt der beste erwartungstreue Schätzer beschreiben.

    \uncover<2->{ $\rightarrow$ Stochastik-Vorlesung}

    \uncover<3->{\vspace{1cm}

      Für viele Verteilungen (z.B. kontaminierte Normalverteilung) gibt es keine geschlossene Formel, nur Iterationsverfahren.

      $\rightarrow$ Maximum-Likelihood-Prinzip
}
  \end{frame}

  
\begin{frame}[separator]
  \frametitle{Vielen Dank für Ihre Aufmerksamkeit!}
\end{frame}

\end{document}

%%% Local Variables:
%%% mode: latex
%%% TeX-master: t
%%% End:
